% Artículo del póster (páginas 1 a 3), formato Springer LNCS/LNAI.
% La página 4 es el póster: ver entrega_4_paginas.tex.
% Todas las cifras salen de Repositorio_Github/results/tables (pipeline src/ueba_ts).
\documentclass[runningheads,a4paper]{llncs}
\usepackage[T1]{fontenc}
\usepackage[utf8]{inputenc}
\usepackage[spanish,es-tabla,es-nodecimaldot]{babel}
\usepackage{graphicx}
\usepackage{booktabs}
\usepackage{amsmath}
\usepackage[hidelinks]{hyperref}
\usepackage{microtype}
\setlength{\textfloatsep}{10pt plus 2pt minus 2pt}
\graphicspath{{figuras/}}
\renewcommand{\keywordname}{{\bfseries Palabras clave:}}

\hypersetup{pdftitle={Detectando lo que las reglas y las líneas base no ven: UEBA basado en series de tiempo},pdfauthor={Alejandro Higuera Castro}}
\begin{document}

\title{\texorpdfstring{\resizebox{\textwidth}{!}{Detectando lo que las reglas y las líneas base no ven:}\\ UEBA basado en series de tiempo}{Detectando lo que las reglas y las líneas base no ven: UEBA basado en series de tiempo}}
\titlerunning{UEBA basado en series de tiempo}
\author{Alejandro Higuera Castro}
\authorrunning{A. Higuera Castro}
\institute{Departamento de Estadística, Facultad de Ciencias\\
Universidad Nacional de Colombia, Sede Bogotá, Colombia\\
\email{ahiguerac@unal.edu.co}}

\maketitle

\begin{abstract}
Los centros de operaciones de seguridad (SOC) buscan amenazas internas con reglas de umbral y líneas base estáticas que generan tantas falsas alarmas que desgastan y desensibilizan a los analistas. Este trabajo trata la actividad de cada usuario como una serie de tiempo con ciclo semanal y aplica la metodología Box-Jenkins a la analítica del comportamiento de usuarios (UEBA): cada noche, un modelo SARIMA pronostica cada uno de los nueve contadores de actividad de cada usuario, mide cuánto se aleja lo observado de lo normal para esa persona y entrega al SOC los casos más anómalos del día. En CERT r4.2 (1.000 usuarios, 500 días, 70 incidentes), con las mismas 10 alertas diarias que una regla global y una línea base estática, SARIMA encuentra 240 días maliciosos, frente a 163 de la estática (47~\% más) y 36 de la regla global, y detecta 67 de los 70 incidentes (estática, 64; global, 17). Tiene potencial para reducir el trabajo de los equipos de seguridad y la fatiga de alertas: la línea base estática necesita 18 alertas diarias para encontrar lo que SARIMA encuentra con 10, y el analista revisa 14 alertas por cada día malicioso encontrado, en vez de 21.
\keywords{Amenaza interna \and UEBA \and Series de tiempo \and SARIMA \and Detección de anomalías \and Fatiga de alertas}
\end{abstract}

\section{Introducción}

Las firmas detectan amenazas conocidas, pero en gran medida no las desconocidas~\cite{nist80094}, como una persona que usa credenciales válidas con fines maliciosos (MITRE ATT\&CK T1078). Los SOC usan reglas de umbral y perfiles~\cite{nist80094}, pero un umbral fijo trata igual un lunes que un domingo: dispara con actividad normal y calla ante anomalías que no lo alcanzan. Casi todas las alarmas son falsos positivos, revisarlas desgasta a los analistas~\cite{alahmadi2022} y así se pierden ataques~\cite{hassan2019nodoze}.

La actividad de un usuario, como el tráfico de red~\cite{brutlag2000}, es una serie de tiempo con ciclo semanal, y su línea base debería modelarla. ¿Con las alertas que un SOC puede revisar, un modelo por usuario detecta más que una regla global y una línea base estática? El aporte son modelos Box-Jenkins interpretables, uno por usuario y contador, evaluados con igual carga de trabajo y con la proporción real de casos (0,32~\% de días maliciosos); no se comparan con redes profundas como BRITD~\cite{song2024britd} y BTITD~\cite{zhang2026btitd}.

\section{Datos y método}

\textbf{Datos.} CERT r4.2~\cite{lindauer2020} simula 1.000 usuarios durante 500 días y registra 70 incidentes: fuga de información a WikiLeaks fuera de horario (30), robo de datos antes de renunciar (30) y sabotaje con la cuenta del supervisor (10). Un usuario-día es malicioso si tiene un evento del incidente.

\textbf{Cómo funciona.} Cada noche: (1) se resumen los registros que ya centraliza el SIEM en nueve contadores por usuario y día (sesiones, USB, archivos, correo y web, también fuera de 07:00 a 19:00); (2) para cada usuario y contador, un modelo pronostica el valor esperado del día y su \emph{límite del 99~\%}, el valor que un día normal no debería superar más de 1 vez de cada 100; (3) el puntaje del usuario es su mayor aumento sobre el pronóstico de log(1 + conteo), en desviaciones estándar; (4) el SOC revisa los $k$ usuarios con mayor puntaje.

\textbf{Modelos.} Se comparan la regla global (un umbral para todos), la línea base estática (media y desviación de cada usuario, recalculadas cada mes), la regla ingenua (hoy igual que hace una semana), ARIMA, SARIMA y SARIMAX (con la mediana del departamento). Según Box-Jenkins, de 5.700 series con actividad, el 94~\% tiene ciclo semanal y en el 96~\% ADF (sin el promedio de cada día) y KPSS son compatibles con estacionariedad; por AIC se eligió SARIMA$(0,0,0)(1,0,1)_7$. Sus coeficientes ($\Phi\approx 1$, $\Theta\approx -1$) lo vuelven un perfil por día de la semana que se actualiza lentamente; las series casi vacías (37~\% de los ajustes) usan la media del usuario.

\textbf{Evaluación.} Se entrena hasta el 2010-06-09, antes de cualquier incidente, y se prueba hasta el 2011-05-16, reentrenando cada mes con 182 días, sin usar etiquetas~\cite{tashman2000}, o, como control, una sola vez. Se mide cuánto se detecta con $k$ alertas diarias y la precisión promedio (PR-AUC; al azar, 0,0032), con intervalos de confianza (IC) del 95~\% por bootstrap de usuarios.

\section{Resultados}

Con igual carga de trabajo, los modelos de series de tiempo superan a las reglas en toda la Tabla~\ref{tab:abl}; con 10 alertas diarias, SARIMA encuentra 240 días maliciosos, frente a 163 de la estática (+47~\%; IC del 95~\% de la diferencia: 43 a 117) y 36 de la regla global. SARIMA supera a la estática con 1 a 30 alertas diarias, aun sin reentrenar, y en PR-AUC (+0,0062; IC: 0,0025 a 0,0108).

\begin{table}[t]
\caption{Detecciones con igual número de alertas diarias. Abajo, enfoques temporales.}
\label{tab:abl}
\centering\scriptsize
\setlength{\tabcolsep}{3.6pt}
\begin{tabular}{lccccccc}
\toprule
 & \multicolumn{4}{c}{Días maliciosos (de 986)} & \multicolumn{2}{c}{Con 10 alertas diarias} & \\
\cmidrule(lr){2-5}\cmidrule(lr){6-7}
Detector & 3/día & 5/día & 10/día & 20/día & Incidentes & Aciertos/100 & PR-AUC \\
\midrule
Regla global & 16 & 19 & 36 & 42 & 17/70 & 1,1 & 0,0125 \\
Línea base estática & 52 & 82 & 163 & 258 & 64/70 & 4,8 & 0,0270 \\
Regla ingenua & 32 & 53 & 93 & 157 & 46/70 & 2,7 & 0,0118 \\
\midrule
Perfil semanal (control) & 75 & 128 & 210 & \textbf{322} & \textbf{67/70} & 6,2 & \textbf{0,0361} \\
ARIMA & \textbf{85} & 126 & 220 & 293 & \textbf{67/70} & 6,5 & 0,0346 \\
SARIMA & 81 & \textbf{146} & \textbf{240} & 300 & \textbf{67/70} & \textbf{7,0} & 0,0331 \\
SARIMAX + pares & 83 & 131 & 234 & 305 & 65/70 & 6,9 & 0,0330 \\
\bottomrule
\end{tabular}
\end{table}

Sin tope diario, la regla global dispara 108 alertas al día y SARIMA, 18,5 (falsas: 98 y 95 de cada 100). La estática necesita 18 alertas diarias para alcanzar los 240 días de SARIMA con 10. El perfil semanal (media y desviación de cada usuario por día de la semana), un control temporal más simple, también supera a la estática con cualquier número de alertas (con 10, 47 días más; IC: 24 a 74); frente a SARIMA encuentra menos con 3 a 17 alertas diarias y más con 18 o más. La ventaja se concentra en el robo de datos antes de renunciar (861 de los 986 días maliciosos): SARIMA encuentra 167 y la estática, 86; en los otros dos escenarios, 73 y 77. El límite es conservador: los días normales lo superaron del 0,2~\% al 0,8~\% de las veces (nominal, 1~\%). En PR-AUC no hubo evidencia de diferencias entre los tres modelos, pero sin reentrenar SARIMA conserva su PR-AUC (0,0339) y ARIMA cae a 0,0204, bajo la estática (0,0277). Ljung-Box no rechaza residuos sin autocorrelación en el 95~\% de los ajustes iniciales de SARIMA y el 82~\% de los mensuales (nominal, 95~\%); en ARIMA, el 9~\%.

\section{Conclusiones}

Un modelo de series de tiempo aprende el ritmo de cada usuario; un umbral fijo, no. Con 10 alertas diarias, los tres modelos detectan entre 35~\% y 47~\% más días maliciosos que la estática, y SARIMA, 6,7 veces los de la regla global; el analista revisa 14,2 alertas por día malicioso encontrado en vez de 20,9, lo que ayuda a reducir la fatiga de alertas. El ritmo semanal hay que estimarlo: copiar la semana anterior encuentra 93 días. No ve costumbres del usuario ni acciones sin cambio de volumen (para eso están la regla global, el EDR y las firmas), y se probó en datos sintéticos: una alerta no confirma un ataque.

\bibliographystyle{splncs04}
\bibliography{referencias}
\end{document}
